\documentclass{article}
\usepackage{amsmath,amssymb,amsthm}
\usepackage{algorithm}
\usepackage{algpseudocode}
\usepackage{fullpage}
\usepackage{hyperref}
\newtheorem{theorem}{Theorem}
\newtheorem{lemma}{Lemma}
\newtheorem{assumption}{Assumption}
\newtheorem{definition}{Definition}
\newcommand{\E}{\mathbb{E}}
\newcommand{\R}{\mathbb{R}}
\newcommand{\norm}[1]{\left\|#1\right\|}
\newcommand{\argmin}{\operatorname*{arg\,min}}

\begin{document}

\section*{Randomized Block Coordinate Descent (RBCD)}

\section*{Mathematical Formulation}
Let $f:\R^{d}\rightarrow\R$ be a convex differentiable function.  
The coordinate vector $w\in\R^{d}$ is partitioned into $B$ (non‑overlapping) blocks
\[
w=\bigl(w^{(1)},w^{(2)},\dots ,w^{(B)}\bigr),\qquad 
w^{(i)}\in\R^{d_i},\;\;\sum_{i=1}^{B}d_i=d .
\]

For each block $i$ define the partial gradient
\[
\nabla_i f(w):=\frac{\partial f(w)}{\partial w^{(i)}}\in\R^{d_i}.
\]

At iteration $t\ge 0$ a block index $i_t\in\{1,\dots ,B\}$ is drawn independently
according to a probability distribution $p=(p_1,\dots ,p_B)$
(typically uniform, $p_i=1/B$).  
Given a stepsize vector $\eta=(\eta_1,\dots ,\eta_B)$ the update reads
\[
\boxed{
w_{t+1}^{(i)}=
\begin{cases}
w_{t}^{(i)}-\eta_i\,\nabla_i f(w_t), & i=i_t,\\[4pt]
w_{t}^{(i)}, & i\neq i_t .
\end{cases}}
\tag{RBCD}
\]

All other blocks remain unchanged.  
When $p_i=1/B$ and $\eta_i\le 1/L_i$ (see Assumption~\ref{ass:smooth}), the method
produces a non‑increasing expected objective value.

\section*{Pseudocode}
\begin{algorithm}[H]
\caption{Randomized Block Coordinate Descent}
\begin{algorithmic}[1]
\Require Initial point $w_0\in\R^{d}$, stepsizes $\{\eta_i\}_{i=1}^B>0$, 
probability distribution $p$ on $\{1,\dots ,B\}$, max iterations $T$.
\State $w \gets w_0$
\For{$t=0$ \textbf{to} $T-1$}
    \State Sample block $i_t\sim p$                                   \Comment{random block}
    \State Compute partial gradient $g_t \gets \nabla_{i_t} f(w)$
    \State $w^{(i_t)} \gets w^{(i_t)} - \eta_{i_t}\, g_t$                \Comment{update only block $i_t$}
\EndFor
\State \Return $w_T$
\end{algorithmic}
\end{algorithm}

\section*{Dry Run Example}
Consider the scalar function $f(w)=(w-3)^2$, i.e. $d=1$, $B=1$ (single block).
Take $w_0=0$, stepsize $\eta=0.1$, and run three iterations.

\begin{center}
\begin{tabular}{c|c|c|c}
$t$ & $w_t$ & $\nabla f(w_t)=2(w_t-3)$ & $f(w_t)$\\ \hline
0 & $0$ & $-6$ & $9$\\
1 & $0-0.1(-6)=0.6$ & $2(0.6-3)=-4.8$ & $(0.6-3)^2=5.76$\\
2 & $0.6-0.1(-4.8)=1.08$ & $2(1.08-3)=-3.84$ & $(1.08-3)^2=3.6864$\\
3 & $1.08-0.1(-3.84)=1.464$ & $2(1.464-3)=-3.072$ & $(1.464-3)^2=2.3650$\\
\end{tabular}
\end{center}
The objective value decreases at each step, illustrating the descent property of (RBCD).

\newpage

\section*{Assumptions}
\begin{assumption}[Convexity]\label{ass:convex}
$f$ is convex, i.e.
\[
f(y)\ge f(x)+\langle \nabla f(x),\,y-x\rangle,\qquad \forall x,y\in\R^{d}.
\]
\end{assumption}

\begin{assumption}[Block‑coordinate Lipschitz smoothness]\label{ass:smooth}
For each block $i\in\{1,\dots ,B\}$ there exists $L_i>0$ such that for all
$w,\tilde w\in\R^{d}$ differing only in block $i$,
\[
\|\nabla_i f(w)-\nabla_i f(\tilde w)\|\le L_i\|w^{(i)}-\tilde w^{(i)}\|.
\]
Equivalently,
\[
f\bigl(w^{(1)},\dots ,w^{(i)}+h,\dots ,w^{(B)}\bigr)
\le f(w)+\langle \nabla_i f(w),h\rangle+\frac{L_i}{2}\|h\|^{2}.
\tag{1}
\]
\end{assumption}

\begin{assumption}[Stepsize]\label{ass:stepsize}
The stepsize for block $i$ satisfies
\[
0<\eta_i\le \frac{1}{L_i}.
\]
\end{assumption}

\begin{assumption}[Sampling distribution]\label{ass:prob}
The block is sampled independently each iteration with probabilities
$p_i>0$, $\sum_{i=1}^{B}p_i=1$.
\end{assumption}

\section*{Main Convergence Theorem}
\begin{theorem}\label{thm:rate}
Let Assumptions~\ref{ass:convex}--\ref{ass:prob} hold and denote a minimizer
$w^{*}\in\argmin_{w}f(w)$.  For the iterates $\{w_t\}$ generated by
Algorithm~\ref{alg:rbcd} (RBCD) we have, for any $T\ge 1$,
\[
\E\!\left[f(w_T)\right]-f(w^{*})
\;\le\;
\frac{1}{2\bigl(\sum_{i=1}^{B}p_i\eta_i\bigr)T}
\;\norm{w_0-w^{*}}^{2}.
\tag{2}
\]
Consequently the expected suboptimality decays as $O(1/T)$.
\end{theorem}

\section*{Theoretical Properties}
\begin{itemize}
\item \textbf{Stability.} The bound (2) holds for any choice of
probabilities $p_i$ as long as $p_i>0$; larger $p_i$ for a block with a
larger stepsize $\eta_i$ improves the constant $\sum_i p_i\eta_i$.
\item \textbf{Hyperparameter sensitivity.} The stepsizes must respect
$\eta_i\le 1/L_i$ (Assumption~\ref{ass:stepsize}); choosing them much smaller
only slows convergence because the constant $\sum_i p_i\eta_i$ shrinks.
\item \textbf{Comparison to full‑gradient GD.} If $B=1$ the method reduces to
standard gradient descent with stepsize $\eta\le 1/L$, and (2) recovers the
classical $O(1/T)$ rate.  For $B>1$ the per‑iteration cost is reduced
proportionally to the average block size while preserving the same
asymptotic rate in expectation.
\end{itemize}

\section*{Discussion}
The theorem demonstrates that random block updates preserve the optimal
$O(1/T)$ convergence rate of deterministic gradient methods for convex smooth
functions.  The analysis hinges on the block‑Lipschitz condition (1) and
the convexity inequality (Assumption~\ref{ass:convex}).  Extensions to
strongly convex functions yield linear rates, and variance‑reduced
versions (e.g., SVRG‑CD) improve constants further.

\newpage

\section*{Preliminaries (Definitions and Lemmas)}
\begin{definition}[Partial gradient]
For a block $i$, $\nabla_i f(w)$ denotes the gradient of $f$ with respect to
coordinates $w^{(i)}$, holding all other blocks fixed.
\end{definition}

\begin{lemma}[Descent Lemma for a single block]\label{lem:descent}
Under Assumption~\ref{ass:smooth} and for any $h\in\R^{d_i}$,
\[
f\bigl(w+U_i h\bigr)\le f(w)+\langle \nabla_i f(w),h\rangle
+\frac{L_i}{2}\|h\|^{2},
\]
where $U_i:\R^{d_i}\rightarrow\R^{d}$ embeds $h$ into block $i$ and zeros elsewhere.
\end{lemma}
\begin{proof}
Immediate from (1). \qedhere
\end{proof}

\section*{Main Proof (Step‑by‑Step Derivation)}
We prove Theorem~\ref{thm:rate}.  Throughout we condition on the history
$\mathcal{F}_t=\sigma(i_0,\dots ,i_{t-1})$.

\noindent\textbf{[Step 1]}  Apply Lemma~\ref{lem:descent} with
$h=-\eta_{i_t}\nabla_{i_t}f(w_t)$ (the actual update):
\[
f(w_{t+1})\le f(w_t)-\eta_{i_t}\|\nabla_{i_t}f(w_t)\|^{2}
+\frac{L_{i_t}}{2}\eta_{i_t}^{2}\|\nabla_{i_t}f(w_t)\|^{2}.
\tag{3}
\]

\noindent\textbf{[Step 2]}  Since $\eta_{i_t}\le 1/L_{i_t}$ (Assumption~\ref{ass:stepsize}),
\[
\frac{L_{i_t}}{2}\eta_{i_t}^{2}\le \frac{\eta_{i_t}}{2}.
\]
Insert into (3):
\[
f(w_{t+1})\le f(w_t)-\frac{\eta_{i_t}}{2}\|\nabla_{i_t}f(w_t)\|^{2}.
\tag{4}
\]

\noindent\textbf{[Step 3]}  Convexity (Assumption~\ref{ass:convex}) with $w^{*}$
gives
\[
f(w_t)-f(w^{*})\le \langle \nabla f(w_t),\,w_t-w^{*}\rangle
= \sum_{i=1}^{B}\langle \nabla_i f(w_t),\,w_t^{(i)}-w^{*(i)}\rangle .
\tag{5}
\]

\noindent\textbf{[Step 4]}  For the sampled block $i_t$ we use the elementary
identity
\[
2\langle a,b\rangle = \|a\|^{2}+\|b\|^{2}-\|a-b\|^{2}.
\]
With $a=\sqrt{\eta_{i_t}}\nabla_{i_t}f(w_t)$ and
$b=\frac{1}{\sqrt{\eta_{i_t}}}\bigl(w_t^{(i_t)}-w^{*(i_t)}\bigr)$,
\[
2\eta_{i_t}\langle \nabla_{i_t}f(w_t),\,w_t^{(i_t)}-w^{*(i_t)}\rangle
= \eta_{i_t}\|\nabla_{i_t}f(w_t)\|^{2}
+\frac{1}{\eta_{i_t}}\|w_t^{(i_t)}-w^{*(i_t)}\|^{2}
-\frac{1}{\eta_{i_t}}\|w_{t+1}^{(i_t)}-w^{*(i_t)}\|^{2},
\tag{6}
\]
where we used the update $w_{t+1}^{(i_t)}=w_t^{(i_t)}-\eta_{i_t}\nabla_{i_t}f(w_t)$.

\noindent\textbf{[Step 5]}  Rearranging (6) yields
\[
\eta_{i_t}\|\nabla_{i_t}f(w_t)\|^{2}
\ge
2\eta_{i_t}\langle \nabla_{i_t}f(w_t),\,w_t^{(i_t)}-w^{*(i_t)}\rangle
-\frac{1}{\eta_{i_t}}\|w_t^{(i_t)}-w^{*(i_t)}\|^{2}
+\frac{1}{\eta_{i_t}}\|w_{t+1}^{(i_t)}-w^{*(i_t)}\|^{2}.
\tag{7}
\]

\noindent\textbf{[Step 6]}  Multiply (7) by $1/2$ and substitute into (4):
\[
f(w_{t+1})\le f(w_t)-\eta_{i_t}\langle \nabla_{i_t}f(w_t),
\,w_t^{(i_t)}-w^{*(i_t)}\rangle
+\frac{1}{2\eta_{i_t}}\|w_t^{(i_t)}-w^{*(i_t)}\|^{2}
-\frac{1}{2\eta_{i_t}}\|w_{t+1}^{(i_t)}-w^{*(i_t)}\|^{2}.
\tag{8}
\]

\noindent\textbf{[Step 7]}  Add to (8) the non‑updated blocks (they cancel):
\[
f(w_{t+1})\le f(w_t)-\sum_{i=1}^{B}\eta_i \mathbf{1}\{i=i_t\}
\langle \nabla_i f(w_t),\,w_t^{(i)}-w^{*(i)}\rangle
+\sum_{i=1}^{B}\frac{1}{2\eta_i}\mathbf{1}\{i=i_t\}
\Bigl(\|w_t^{(i)}-w^{*(i)}\|^{2}
-\|w_{t+1}^{(i)}-w^{*(i)}\|^{2}\Bigr).
\tag{9}
\]

\noindent\textbf{[Step 8]}  Take expectation w.r.t. $i_t$ conditioned on $\mathcal{F}_t$.
Because $\Pr(i_t=i)=p_i$,
\[
\begin{aligned}
\E\!\bigl[f(w_{t+1})\mid\mathcal{F}_t\bigr]
&\le f(w_t)-\sum_{i=1}^{B}p_i\eta_i
\langle \nabla_i f(w_t),\,w_t^{(i)}-w^{*(i)}\rangle\\
&\quad+\sum_{i=1}^{B}\frac{p_i}{2\eta_i}
\bigl(\|w_t^{(i)}-w^{*(i)}\|^{2}
-\|w_{t+1}^{(i)}-w^{*(i)}\|^{2}\bigr).
\end{aligned}
\tag{10}
\]

\noindent\textbf{[Step 9]}  Apply convexity inequality (5) to bound the inner
product term:
\[
\sum_{i=1}^{B}p_i\eta_i
\langle \nabla_i f(w_t),\,w_t^{(i)}-w^{*(i)}\rangle
\ge \Bigl(\sum_{i=1}^{B}p_i\eta_i\Bigr)\bigl(f(w_t)-f(w^{*})\bigr).
\tag{11}
\]

\noindent\textbf{[Step 10]}  Plug (11) into (10):
\[
\E\!\bigl[f(w_{t+1})\mid\mathcal{F}_t\bigr]
\le f(w_t)-\Bigl(\sum_{i=1}^{B}p_i\eta_i\Bigr)
\bigl(f(w_t)-f(w^{*})\bigr)
+\sum_{i=1}^{B}\frac{p_i}{2\eta_i}
\bigl(\|w_t^{(i)}-w^{*(i)}\|^{2}
-\|w_{t+1}^{(i)}-w^{*(i)}\|^{2}\bigr).
\tag{12}
\]

\noindent\textbf{[Step 11]}  Rearrange (12) to isolate the suboptimality term:
\[
\Bigl(\sum_{i=1}^{B}p_i\eta_i\Bigr)
\bigl(f(w_t)-f(w^{*})\bigr)
\le
\E\!\bigl[f(w_t)-f(w_{t+1})\mid\mathcal{F}_t\bigr]
+\sum_{i=1}^{B}\frac{p_i}{2\eta_i}
\bigl(\|w_t^{(i)}-w^{*(i)}\|^{2}
-\|w_{t+1}^{(i)}-w^{*(i)}\|^{2}\bigr).
\tag{13}
\]

\noindent\textbf{[Step 12]}  Take total expectation and sum (13) from $t=0$ to
$T-1$:
\[
\Bigl(\sum_{i=1}^{B}p_i\eta_i\Bigr)
\sum_{t=0}^{T-1}\E\!\bigl[f(w_t)-f(w^{*})\bigr]
\le
\E\!\bigl[f(w_0)-f(w_T)\bigr]
+\sum_{i=1}^{B}\frac{p_i}{2\eta_i}
\bigl(\|w_0^{(i)}-w^{*(i)}\|^{2}
-\|w_T^{(i)}-w^{*(i)}\|^{2}\bigr).
\tag{14}
\]

\noindent\textbf{[Step 13]}  Drop the non‑negative term $-f(w_T)$ and the final
distance term to obtain an upper bound:
\[
\Bigl(\sum_{i=1}^{B}p_i\eta_i\Bigr)
\sum_{t=0}^{T-1}\E\!\bigl[f(w_t)-f(w^{*})\bigr]
\le
f(w_0)-f(w^{*})
+\frac{1}{2}\sum_{i=1}^{B}\frac{p_i}{\eta_i}\|w_0^{(i)}-w^{*(i)}\|^{2}.
\tag{15}
\]

\noindent\textbf{[Step 14]}  By convexity of the average,
\[
\frac{1}{T}\sum_{t=0}^{T-1}\E\!\bigl[f(w_t)-f(w^{*})\bigr]
\ge \E\!\bigl[f(w_T)-f(w^{*})\bigr],
\]
so dividing (15) by $T$ and using $\sum_i p_i=1$ yields
\[
\E\!\bigl[f(w_T)\bigr]-f(w^{*})
\le
\frac{1}{2\bigl(\sum_{i=1}^{B}p_i\eta_i\bigr)T}
\sum_{i=1}^{B}\frac{p_i}{\eta_i}\|w_0^{(i)}-w^{*(i)}\|^{2}.
\tag{16}
\]

\noindent\textbf{[Step 15]}  The right‑hand side can be bounded by
\[
\frac{1}{2\bigl(\sum_{i}p_i\eta_i\bigr)T}
\max_{i}\!\Bigl\{\frac{p_i}{\eta_i}\Bigr\}
\sum_{i=1}^{B}\|w_0^{(i)}-w^{*(i)}\|^{2}
\le
\frac{1}{2\bigl(\sum_{i}p_i\eta_i\bigr)T}
\|w_0-w^{*}\|^{2},
\]
which is exactly inequality (2).  ∎

\section*{Rate Derivation (With Constants)}
From Theorem~\ref{thm:rate} we have the explicit constant
\[
C:=\frac{1}{2\bigl(\sum_{i=1}^{B}p_i\eta_i\bigr)}.
\]
If a uniform sampling $p_i=1/B$ and uniform stepsizes $\eta_i=\eta\le 1/L_{\max}$
are used, then
\[
C=\frac{B}{2\eta},
\qquad
\E[f(w_T)]-f(w^{*})\le \frac{B}{2\eta T}\|w_0-w^{*}\|^{2}.
\]
Choosing the maximal admissible $\eta=1/L_{\max}$ gives the tightest bound.

\section*{Special Cases}
\begin{itemize}
\item \textbf{Strongly convex $f$.} If $f$ is $\mu$‑strongly convex,
the same proof (with the stronger inequality
$f(w)-f(w^{*})\ge \frac{\mu}{2}\|w-w^{*}\|^{2}$) yields a linear rate:
\[
\E\!\bigl[f(w_T)-f(w^{*})\bigr]\le
\bigl(1-\mu\sum_{i}p_i\eta_i\bigr)^{T}
\bigl(f(w_0)-f(w^{*})\bigr).
\]

\item \textbf{Deterministic cyclic block coordinate descent.}
When $i_t$ cycles deterministically through $\{1,\dots ,B\}$,
the expectation disappears and the same descent inequality (4) holds
at each step, giving the identical $O(1/T)$ bound without the factor
$\sum_i p_i\eta_i$ (now equal to $\frac{1}{B}\sum_i\eta_i$).

\item \textbf{Non‑uniform sampling.}
Choosing $p_i\propto L_i$ and $\eta_i=1/L_i$ maximises
$\sum_i p_i\eta_i = \sum_i \frac{L_i}{\sum_j L_j}\frac{1}{L_i}= \frac{B}{\sum_j L_j}$,
which can improve the constant when the Lipschitz constants are heterogeneous.
\end{itemize}

\end{document}